\clearpage\begingroup\SourceLeaf{47b}\HistoricalHead{42}{L}{PHYSIK. BRIEFWECHSEL.}
\begin{center}
[7.]\\[.7em]
\textsc{Gerling} an \textsc{Gauss}. Marburg, 28. Januar 1853.
\end{center}
\noindent .\ .\ .\ .\ .\ .\ .\ .\ .\ .\ .\ .\ .\ .\ .\ .\ .\ .\ .\ .\ .
\par\noindent\hspace*{\parindent}\{Sehr hat es mich erfreut zu lesen, dass Sie einen Apparat für die \textsc{Foucault}schen
Versuche wollen ausführen lassen. Hoffentlich wird dann auch
mehr Ernst und Genauigkeit in die Versuche selbst kommen, als mir bisher
dabei angewandt zu seyn scheint. Was ich für diesen Zweck früher, und wie
ich mir schmeichle gut, vorbereitet hatte, hat nun schon über ein Jahr geruht,
da die beiden Bedingungen der wirklichen Anwendung, ein ständiger Gehülfe,
den ich zu jeder Zeit zur Assistenz herbeicitiren könnte, und die eigene,
seit längerer Zeit mehr oder weniger gestörte, im letzten Jahr aber fast gänzlich
absorbirte Geistesruhe, die zu soliden Arbeiten unumgänglich nöthig ist,
mir fehlten.

Vielen Dank sage ich Ihnen für die ausführliche Auseinandersetzung,
weshalb sich die Theorie dieser Versuche nicht in ein Aperçu bringen lässt. —
Sie beantworten eigentlich meine Frage: »Was ist dabei als constant vorauszusetzen?«
dahin: »Es ist nichts constant zu setzen als die Anziehung nach
\clearpage\SourceLeaf{48}\HistoricalHead{43}{R}{GERLING AN GAUSS.}
\noindent dem Mittelpunkte der Erde und die Entfernung von dem im Kreise gleichförmig
umlaufenden Aufhängepunkt«. Die Entstehung meiner Frage schreibt
sich eigentlich aus einer ziemlichen Anzahl von Auseinandersetzungen her, die
mehr oder weniger Aehnlichkeit mit der von Ihnen gerügten \textsc{Schaub}schen in
No. 838 [der Astronomischen Nachrichten] darbieten. Dieselben gehen denn
darauf aus, dass bewiesen werden soll, der Winkel, den die successiven Lagen
der Mittagslinien durchlaufen, verhalte sich zu dem Winkel, den die correspondirenden
successiven Lagen der Meridianebene machen, wie \(\sin\varphi:1\).
Da dies aber ein bloss geometrischer Elementarsatz ist, so fragte ich: wie
das mit dem \textls[120]{Pendel} zusammenhänge, und weil die specifische Verschiedenheit
von dem Fall \(\varphi=90^0\) mir nicht klar war, so drückte ich das so aus:
»Für \(\varphi=90^0\) sagt ihr, die Pendel-Ebene sey constant, was ist aber nun constant
zu setzen, wenn \(\varphi<90^0\)? Ehe ihr mir dies erklärt, kann ich nicht
einsehen, dass jener geometrische Satz für die Aufgabe genügt«. — Auf \textsc{Hansens}
wie ich sehe inzwischen in Danzig gekrönte Abhandlung[\textsuperscript{*})]\begingroup\renewcommand{\thefootnote}{}\footnotetext{[\textsuperscript{*}) P. A. \textsc{Hansen}, \textit{Theorie der Pendelbewegung mit Rücksicht auf die Gestalt und Bewegung der Erde}, Neue Schriften der Naturforschenden Gesellschaft zu Danzig \textbf{5}, Heft 1, 1856; auch in \textsc{Poggendorffs} Annalen der Physik und Chemie \textbf{92}, 1854, S. 21.]}\endgroup\ bin ich sehr
neugierig. Er hat im Herbst 1851 in Gotha mit meinem Kollegen \textsc{Stegmann},
wie ich höre, über den Gegenstand gestritten; und letzterer konnte mir damals,
als ich ihm die obige Frage einwarf, auch keine Auskunft geben. H.
wird also wohl den von Ihnen angegebenen specifischen Unterschied gehörig
herausheben. — \textsc{Clausens} Abhandlung habe ich bis jetzt nicht gesehen, werde
aber, sobald ich wieder ausgehen kann, einmal auf der Bibliothek danach
suchen.

Der andere Gegenstand, über den Sie mich belehren, verpflichtet mich
zu nicht minderem Dank, obwohl ich beschämt fühle, dass ich bei grösserer
Aufmerksamkeit (namentlich auf die Schlebuscher Versuche) doch vielleicht
noch selbst dazu gekommen seyn würde einzusehen, wie wenig Grund vorhanden
sey, in die bewussten Scrupel einzustimmen. Dies ist mir besonders
klar geworden, als ich in Folge Ihres Briefes die betreffenden Stellen in \textsc{Benzenberg}
und \textsc{Reich} wieder nachschlug. — Ganz überraschend war aber der Umstand,
dass \textsc{Reichs} Rechnungs-Resultat sogar falsch befunden wurde. — Ich
stelle mir vor, dass es von Nutzen seyn würde, wenn Sie auch dieses gelegent-\HistoricalFoot{}{6\textsuperscript{*}}
\clearpage\SourceLeaf{49}\HistoricalHead{44}{L}{PHYSIK. BRIEFWECHSEL.}
\noindent lich veröffentlichen liessen, denn es ist gar leicht möglich, dass Viele, die mit
mir auf denselben falschen Weg geraten seyn können, dadurch gewarnt werden.
— Bei der Gelegenheit erlaube ich mir noch die Frage, ob es auf einem
Fehler meines Gedächtnisses (das leider sehr abgenommen hat) beruht, wenn
ich meine, vor ca. 8 Jahren gelesen zu haben, dass auch \textsc{Baily} diese Versuche
wiederholt oder ihre Wiederholung veranlasst habe? — Die Wiederholung der
anderen Versuche von \textsc{Reich} über Dichtigkeit der Erde durch \textsc{Baily} ist mir
ganz deutlich in der Erinnerung. — Jenes aber schwebt mir nur ganz dunkel
vor.\quad .\ .\ .\ .\ .\ .\ .\ \hspace{.6em}.\ .\ .\ .\ .\ .\ .\ .\ .\ .\ .\ .\ .\}
\clearpage\endgroup
